% Source presentations for register row N-c-013.

\subsection{The automata-theoretic verification architecture}
\label{sec:p09-presentation-n-c-013}

\subsubsection{Vardi and Wolper's 1986 presentation}

\paragraph{Source state and locators.}
The IEEE CEDA record links an eighteen-page full-text file, while Wolper's
University of Li\`ege repository identifies the same file as an author
preprint.  Its abstract cover is dated 15 August 1985; the separate
proceedings record is the 1986 publication on pp.~322--331.  The accessible
file is therefore not described as a facsimile of those ten printed pages.
The locators below are the preprint's own body and appendix pages
\cite{VardiWolper1986}.

\paragraph{The source language and acceptor.}
 Let \(\mathit{Prop}\) be a finite set of atomic propositions; for a fixed
 formula it may be taken to contain exactly the letters occurring in that
 formula.  A computation is a
function
\[
 \Pi:\omega\longrightarrow 2^{\mathit{Prop}}.
\]
The source's propositional temporal language is closed under Boolean
operations, next \(X\), and until \(U\).  In addition to the ordinary Boolean
clauses, Appendix~A prints
\[
\begin{aligned}
 \Pi,i\models X\psi
 &\Longleftrightarrow \Pi,i+1\models\psi,\\
 \Pi,i\models \xi U\psi
 &\Longleftrightarrow
   \exists j\ge i\,
   \bigl(\Pi,j\models\psi\land
         \forall k\,(i\le k<j\Rightarrow\Pi,k\models\xi)\bigr).
\end{aligned}
\]
Thus the right operand of until may already hold at the current point, and
\(\Pi\models\varphi\) means \(\Pi,0\models\varphi\).  This is the source
presentation of the non-strict infinite-word semantics reused in
Definition~\ref{def:p09-ltl}; it is not Kamp's strict order semantics.

The source B\"uchi automaton is
\[
 A=(\Sigma,S,\rho,S_0,F),
 \qquad
 \rho:S\times\Sigma\longrightarrow 2^S,
\]
with a set \(S_0\subseteq S\) of starting states.  A run on
\(a_1a_2\cdots\) is \(s_0s_1\cdots\), where
\(s_0\in S_0\) and \(s_i\in\rho(s_{i-1},a_i)\) for every \(i\ge1\).
It accepts when some state in \(F\) occurs infinitely often.  The source calls
the automaton \emph{semi-deterministic} exactly when
\[
 |\rho(s,a)|\le1
 \quad\text{for every }s\in S\text{ and }a\in\Sigma,
\]
and deterministic when this holds and \(|S_0|=1\).  Semi-determinism stands
as a restriction on the transition data and differs from determinism only by
permitting several starts; Section~4.4 uses precisely that distinction
\cite[preprint pp.~4 and 10 and app.~A]{VardiWolper1986}.

The handbook's automaton convention has one initial state.  For any
permitted-start-set automaton \((Q,\Sigma,\Delta,Q_0,F)\) in the handbook's
zero-based transition convention, choose a fresh nonaccepting state \(q_*\) and
put
\[
 Q_*=Q\cup\{q_*\},
 \qquad
 F_*=F,
 \qquad
 \Delta_*(q,a)=
 \begin{cases}
  \displaystyle\bigcup_{q_0\in Q_0}\Delta(q_0,a),&q=q_*,\\
  \Delta(q,a),&q\in Q.
 \end{cases}
\]
No transition has target \(q_*\).  A run from \(q_*\) consumes the first
letter while choosing and advancing from one permitted source start.  Removing
or inserting the resulting single visit at position zero cannot change which
old states recur infinitely often.  Thus this handbook-owned conversion
preserves the infinite-word language, also when \(Q_0=\varnothing\), and adds
exactly one state.  It will be used both for the Vardi--Wolper product and for
the finite data extracted from B\"uchi's recurrent formulas below.

Theorem~2.1, which the preprint attributes to Wolper--Vardi--Sistla's 1983
paper, states that every temporal formula \(\varphi\) has an effectively
constructed B\"uchi automaton \(A_\varphi\) over
\(2^{\mathit{Prop}}\) such that
\[
 |S|\le2^{O(|\varphi|)},
 \qquad
 L(A_\varphi)=\{\Pi:\Pi\models\varphi\}.
\]
This is the literal source bound; Appendix~B of the preprint supplies the
compact construction below, while the theorem's earlier attribution remains
visible \cite{WolperVardiSistla1983,VardiWolper1986}.  Theorem~2.2 separately cites linear-time
and nondeterministic-logarithmic-space procedures for B\"uchi emptiness
\cite[preprint p.~4]{VardiWolper1986}.

\paragraph{What Appendix B constructs.}
For a formula \(\varphi\), let \(cl(\varphi)\) be its subformulas and their
negations, with double negations identified, and let \(e(\varphi)\) be the
until formulas in that closure.  Put
\[
 B_\varphi=2^{cl(\varphi)}.
\]
This is the intermediate alphabet of complete closure types; it differs from
the final alphabet \(2^{\mathit{Prop}}\), and the projection between them is
defined below.  The appendix records \(|cl(\varphi)|\le2|\varphi|\).

The local-state set \(N_L\) consists exactly of the subsets
\(u\subseteq cl(\varphi)\) satisfying the following scoped conditions:
\[
\begin{aligned}
 &\psi\in u\Longleftrightarrow\neg\psi\notin u
 &&\text{for every }\psi\in cl(\varphi),\\
 &\xi\land\psi\in u
   \Longleftrightarrow \xi\in u\land\psi\in u
 &&\text{for every }\xi\land\psi\in cl(\varphi).
\end{aligned}
\]
Put
\[
 N_\varphi=\{u\in N_L:\varphi\in u\}.
\]
Thus \(N_\varphi\) is the permitted initial subset.  The names \(B_\varphi\),
\(N_L\), and \(N_\varphi\) are eliminable abbreviations for the displayed
sets, not additional logical primitives.

Define the local transition
\[
 \rho_L:N_L\times B_\varphi\longrightarrow\mathcal P(N_L)
\]
by declaring, for \(u,v\in N_L\) and \(a\in B_\varphi\), that
\(v\in\rho_L(u,a)\) exactly when \(a=u\) and
\[
\begin{aligned}
 X\psi\in u
 &\Longleftrightarrow \psi\in v
 &&\text{for every }X\psi\in cl(\varphi),\\
 \xi U\psi\in u
 &\Longleftrightarrow
   \psi\in u\ \lor\
   (\xi\in u\land \xi U\psi\in v)
 &&\text{for every }\xi U\psi\in cl(\varphi).
\end{aligned}
\]
These clauses check Boolean and one-step consistency but do not by themselves
force a pending until formula ever to reach its right operand.

The eventuality component has state set \(2^{e(\varphi)}\), reads
\(B_\varphi\), starts at the empty set, and designates only the empty set.
Its transition
\[
 \rho_E:2^{e(\varphi)}\times B_\varphi
 \longrightarrow\mathcal P(2^{e(\varphi)})
\]
is defined, for current obligation set \(v\), current closure-type letter
\(a\), and next obligation set \(z\), by
\[
 z\in\rho_E(v,a)
 \Longleftrightarrow
 \begin{cases}
  \forall\,\xi U\psi\in e(\varphi)\cap a\;
  (\psi\in a\lor \xi U\psi\in z),&v=\varnothing,\\
  \forall\,\xi U\psi\in v\;
  (\psi\in a\lor \xi U\psi\in z),&v\ne\varnothing.
 \end{cases}
\]
When the current set is empty, the first clause requires every currently
unsatisfied until formula to enter the next batch.  While a batch is nonempty,
the second clause requires each member whose right operand is still false to
remain, but requires no newly encountered eventuality to be added.  Because
the transition is nondeterministic, an accepting run may choose no additional
members.  Requiring the empty set to recur infinitely often therefore
discharges each required batch before the next one is collected.

The product automaton has
\[
 N_M=N_L\times2^{e(\varphi)},
 \qquad
 N_{M0}=N_\varphi\times\{\varnothing\},
 \qquad
 F_M=N_L\times\{\varnothing\},
\]
and transition
\[
 \rho_M((u,v),a)
 =\rho_L(u,a)\times\rho_E(v,a)
 \qquad
 ((u,v)\in N_M,\ a\in B_\varphi).
\]
Its explicit bound is
\[
 |N_M|
 \le2^{|cl(\varphi)|}2^{|e(\varphi)|}
 \le2^{3|\varphi|}.
\]
The last inequality uses \(|e(\varphi)|\le|\varphi|\), since every member of
\(e(\varphi)\) is an until subformula.  This is the source product with the
permitted initial set \(N_\varphi\times\{\varnothing\}\); applying the
handbook-owned single-start conversion above produces \(|N_M|+1\) states and
does not alter the language.
It initially reads the complete closure types in \(B_\varphi\).  Define
\[
 \pi:B_\varphi\longrightarrow2^{\mathit{Prop}},
 \qquad
 \pi(b)=b\cap\mathit{Prop},
\]
and, for \(c\in2^{\mathit{Prop}}\), define the projected transition by
\[
 \rho_M^\pi((u,v),c)
 =\bigcup_{\substack{b\in B_\varphi\\\pi(b)=c}}
   \rho_M((u,v),b).
\]
The projected automaton retains \(N_M,N_{M0},F_M\) and reads propositional
valuations.  Appendix~B supplies this construction and its obligation
invariant, but not a line-by-line induction for both language inclusions or a
worked input formula
\cite[app.~B, pp.~2--3]{VardiWolper1986}.

\paragraph{Finite-state programs and the product theorem.}
A finite-state program in Section~3.1 is
\[
 P=(W,s_0,R,V),
\]
where \(W\) is finite, \(s_0\in W\), \(R\subseteq W^2\) is total, and
\(V:W\to2^{\mathit{Prop}}\).  An infinite \(R\)-path from \(s_0\) becomes a
computation by reading its state labels, and \(P\) satisfies \(\varphi\) when
every such computation satisfies it.  The program itself is the B\"uchi
automaton
\[
 A_P=(2^{\mathit{Prop}},W,s_0,\rho,W),
 \qquad
 s'\in\rho(s,a)
 \Longleftrightarrow
 (s,s')\in R\land a=V(s).
\]
This display puts the singleton start \(s_0\) before \(\rho\), whereas the
general tuple above puts \(\rho\) before the start set \(S_0\).  It is a
source-order inconsistency, not a second automaton convention; the typed data
are alphabet \(2^{\mathit{Prop}}\), states \(W\), transition \(\rho\), start
set \(\{s_0\}\), and designated set \(W\).
Every infinite run is accepting because every program state is designated.

The required language inclusion and its emptiness form are
\[
 L(A_P)\subseteq L(A_\varphi)
 \quad\Longleftrightarrow\quad
 L(A_P)\cap L(A_{\neg\varphi})=\varnothing.
\]
No automata-theoretic complementation is needed: Theorem~2.1 is applied
directly to \(\neg\varphi\).  An accepting intersection run is a legal
program computation satisfying \(\neg\varphi\), hence a counterexample.
Conversely, every counterexample computation is accepted by both factors.
This is the two-direction correctness argument underlying
Definition~\ref{def:p09-ltl-model-checking}.

Lemma~3.1 cites an intersection construction with \(2|S||S'|\) states for
two B\"uchi automata.  Because \(A_P\) designates all its states, the source
observes that \(|S||S'|\) states suffice here.  The resulting automaton has
\(|W|2^{O(|\varphi|)}\) states.  Theorem~3.2 gives the two-parameter bounds
\[
\begin{aligned}
 \text{time}&=O(\lVert P\rVert\,2^{O(|\varphi|)}),\\
 \text{space}&=O\bigl((\log\lVert P\rVert+|\varphi|)^2\bigr),
\end{aligned}
\]
where \(\lVert P\rVert\) is the encoded finite-program size.  This theorem is
an upper bound.  Sistla--Clarke's separate source problem asks whether some
path in a finite \(R\)-structure satisfies the formula, exactly the
counterexample half used here.  Their PSPACE-hardness result supplies the
combined-input lower bound; closure of deterministic polynomial space under
complement transfers it to the universal-path statement in
Theorem~\ref{thm:p09-ltl-model-checking}.  Neither lower bound is silently
attributed to Vardi--Wolper's Theorem~3.2
\cite[preprint p.~5]{VardiWolper1986}.

\paragraph{A complete product calculation.}
The source gives no worked product graph, so the following is a handbook
calculation rather than a source example.  Let
\[
 W=\{s_0,g,b\},
 \quad
 R=\{(s_0,g),(s_0,b),(g,g),(b,b)\},
\]
and let \(p\) hold only at \(g\).  Consider the specification \(Fp\).  A
deterministic B\"uchi automaton for its negation \(G\neg p\) has an accepting
state \(q\) and a nonaccepting sink \(d\): a letter omitting \(p\) leaves
\(q\) at \(q\), a letter containing \(p\) sends it to \(d\), and \(d\)
loops on every letter.  The product path
\[
 (s_0,q),(b,q),(b,q),\ldots
\]
reaches the accepting cycle at \((b,q)\).  Its projected program path stays at
\(b\), so \(p\) never occurs and \(Fp\) fails.  The other branch reaches
\((g,d)\) and cannot accept.  Thus the accepting product cycle is not merely
a graph artifact: it is exactly the ultimately periodic counterexample
\(s_0b^\omega\).

\paragraph{Past operators and automata connectives.}
Section~3.2 adds previous \(Y\) and since \(S\) with the exact clauses already
given in Definition~\ref{def:p09-ltl}.  Appendix~B modifies the local
transition and initial-state clauses, and the source says Theorem~2.1 and the
model-checking bounds continue to hold.

Section~3.3 instead permits a nondeterministic finite automaton to stand as an
\(n\)-ary temporal connective.  The three resulting source languages differ
by the connective automaton's acceptance:
\[
\begin{array}{c@{\quad}l}
 ETL_f & \text{some finite run on a prefix ends in }F,\\
 ETL_l & \text{some infinite run exists on the word},\\
 ETL_r & \text{some infinite run visits }F\text{ infinitely often}.
\end{array}
\]
The names stand for formula languages, not program classes.  The distinction
is necessary because finite success, continued running, and infinite
recurrence are different conditions, and the source uses all three in its
translation comparison.  It says the analogue of Theorem~2.1 was known for
\(ETL_f\) and \(ETL_l\); Theorem~3.3 gives for \(ETL_r\)
\[
 |S|\le O\bigl(\exp(|\varphi|^2)\bigr).
\]
That cited translation proof is not reproduced in the paper
\cite[preprint p.~6 and app.~A]{VardiWolper1986}.

\paragraph{Measure-one satisfaction is a different problem.}
The probabilistic section begins with a labelled Markov chain
\[
 \Pi=(W,P,w_0,V),
 \qquad
 W\text{ finite},
 \qquad
 w_0\in W,
 \qquad
 P:W^2\to[0,1],
 \qquad
 (\forall u\in W)\ \sum_{v\in W}P(u,v)=1,
 \qquad
 V:W\to\Sigma.
\]
Define the path space
\[
 \Omega_\Pi=
 \{\omega\in W^\omega:\omega(0)=w_0\text{ and }
 P(\omega(i),\omega(i+1))>0\text{ for every }i\}.
\]
For every finite sequence \(w_0,\ldots,w_n\) with
\(P(w_i,w_{i+1})>0\) for \(i<n\), its prefix-cylinder event is
\[
 C_\Pi(w_0,\ldots,w_n)
 =\{\omega\in\Omega_\Pi:\omega(i)=w_i\text{ for }0\le i\le n\}.
\]
Let \(\mathcal B_\Pi\) be the Borel \(\sigma\)-algebra on \(\Omega_\Pi\)
generated by these cylinders, and let \(\mu_\Pi\) be the probability measure
determined by
\[
 \mu_\Pi(C_\Pi(w_0,\ldots,w_n))
 =\prod_{i=0}^{n-1}P(w_i,w_{i+1}).
\]
For a probabilistic program \(\Sigma=2^{\mathit{Prop}}\).  Put
\[
 E_\Pi(\varphi)
 =\{\omega\in\Omega_\Pi:
 V(\omega(0))V(\omega(1))\cdots\models\varphi\}.
\]
This is the measurable formula event.  The source's measure-one semantics is
\[
 \Pi\models\varphi
 \Longleftrightarrow
 \mu_\Pi(E_\Pi(\varphi))=1.
\]
This is almost-sure satisfaction.  A measure-zero counterexample set may be
nonempty, so ordinary language emptiness is not substituted for this clause.
The distinct abbreviations \(C_\Pi\) and \(E_\Pi\) replace the source's
overloaded \(\Delta\) notation so that a prefix cylinder is not confused with
a formula event.  In this probabilistic subsection define
\[
 \lVert\Pi\rVert_+
 =|\{(u,v)\in W^2:P(u,v)>0\}|.
\]
This is the source's finite program-size parameter: the algorithms use only
the support of the transition matrix, not the magnitudes of its positive
entries.
These measure and size conventions are on preprint p.~7
\cite[preprint p.~7]{VardiWolper1986}.

A Streett automaton replaces the B\"uchi set by a finite family
\(\mathcal F\subseteq(2^S)^2\).  A run accepts exactly when, for every
\((L,U)\in\mathcal F\), infinite recurrence in \(L\) implies infinite
recurrence in \(U\).  The pair \((S,F)\) represents B\"uchi acceptance.  The
next sentence in the preprint calls a conversion ``opposite'' but its
parenthesis again says B\"uchi-to-Streett while assigning an exponential
cost.  The parenthesis repeats the already trivial direction; the intended
expensive direction can only be Streett-to-B\"uchi.  The handbook records the
printed defect instead of quoting it as a coherent theorem
\cite[preprint p.~8]{VardiWolper1986}.

\paragraph{Restricted probabilistic theorems.}
Put \(F\varphi\equiv \mathit{true}\,U\,\varphi\).  Here \(TL(F)\),
\(TL(X,F)\), and \(TL(Y,S)\) are syntactic formula classes generated from
atomic propositions by Boolean connectives and, respectively, by \(F\) only,
by \(X\) and \(F\) only, or by \(Y\) and \(S\) only.  The class
\(past\text{-}TL(X,F)\) is generated by Boolean connectives and
\(X,F,Y,S\), subject to the condition that no occurrence of \(X\) or \(F\)
lies in the scope of an occurrence of \(Y\) or \(S\).  These names stand for
formula fragments, differ exactly in their permitted temporal constructors,
and are used below to delimit the source's four restricted results.
For an infinite word
\(w\), let \(\operatorname{let}(w)\) be the letters occurring in it and
\(\operatorname{lim}(w)\) those occurring infinitely often.  Two words are
\emph{limit equivalent} exactly when, for some
\(u\in\Sigma^*\), \(a\in\Sigma\), and infinite tails \(v,v'\),
\[
 w=uav,
 \qquad
 w'=uav',
 \qquad
 \operatorname{let}(av)
 =\operatorname{lim}(w)
 =\operatorname{lim}(w')
 =\operatorname{let}(av').
\]
This source relation stands between infinite words; it is weaker than word
equality because it discards order and finite multiplicity in the tails.  Its
distinction is used immediately in Lemma~4.1: limit-equivalent time
structures satisfy the same \(TL(F)\) formulas at zero.

A computation is \emph{probabilistically fair} when every positive-probability
choice out of a state is taken infinitely often whenever that state itself
occurs infinitely often.  This stands as a condition on a computation, not as
universal quantification over paths.  It is stronger than mere recurrence and
is used to make one computation represent a positive-measure class.  Streett
pairs express the condition.

Lemma~4.2 gives a proof sketch: equip the program automaton with that Streett
condition, intersect it with \(A_{\neg\varphi}\), and test Streett emptiness.
Theorem~4.3 concludes for \(TL(F)\)
\[
 \text{time}=O(\lVert\Pi\rVert_+\,2^{O(|\varphi|)}).
\]
The paper explicitly leaves open whether Theorem~3.2's space bound survives
in this probabilistic setting
\cite[preprint pp.~8--9]{VardiWolper1986}.

For \(TL(X,F)\), a formula with \(k-1\) occurrences of \(X\) is evaluated
through length-\(k\) factors.  Write \(k\text{-}\operatorname{let}(w)\) for
the length-\(k\) factors which occur in \(w\), and
\(k\text{-}\operatorname{lim}(w)\) for those which occur infinitely often.
The source defines \(k\)-limit equivalence exactly by the existence of
\(u\in\Sigma^*\), \(x\in\Sigma^k\), and infinite tails \(v,v'\) such that
\[
 w=uxv,
 \quad
 w'=uxv',
 \quad
 k\text{-}\operatorname{let}(xv)
 =k\text{-}\operatorname{lim}(w)
 =k\text{-}\operatorname{lim}(w')
 =k\text{-}\operatorname{let}(xv').
\]
It differs from ordinary limit equivalence by remembering exactly the finite
window needed for the next operators, and Lemma~4.4 uses it to preserve truth
of the formula.  The program is augmented to remember a length-\(k\) history;
details are deferred to the full paper.  Theorem~4.5 states
\[
 \text{time}=O(\lVert\Pi\rVert_+\,2^{O(|\varphi|^2)})
\]
and leaves improvement open.

For the pure-past fragment \(TL(Y,S)\), Lemma~4.6 says that
\(A_\varphi\) is semi-deterministic.  Theorem~4.7 gives
\[
\begin{aligned}
 \text{time}&=O(\lVert\Pi\rVert_+\,2^{O(|\varphi|)}),\\
 \text{space}&=O\bigl((\log\lVert\Pi\rVert_++|\varphi|)^2\bigr).
\end{aligned}
\]
The broader \(past\text{-}TL(X,F)\) fragment forbids \(X\) or \(F\) inside
the scope of \(Y\) or \(S\).  The paper says this fragment has the same
expressive power as PTL but that the cited translation from PTL can have
nonelementary complexity.  Its claimed model-checking extension is deferred
to the full paper, as is the algorithm for the displayed concurrent
Markov-chain model
\[
 (W,N,F,P,w_0,V),
 \qquad
 N\subseteq W,
 \qquad
 F\subseteq N,
 \qquad
 P:W^2\to[0,1],
 \qquad
 w_0\in W,
 \qquad
 V:W\to\Sigma.
\]
For \(u\in W\setminus N\), the outgoing probabilities sum to one.  For
\(u\in N\), the condition \(P(u,v)>0\) records a possible nondeterministic
transition rather than its probability; \(F\) marks the nondeterministic
states subject to fair scheduling
\cite[preprint pp.~10--11]{VardiWolper1986}.

The proof boundary is therefore uneven and explicit.  Theorem~2.1 is
attributed to Wolper--Vardi--Sistla 1983 while Appendix~B constructs the PTL
automaton; Theorems~2.2 and 3.3 and Lemmas~3.1 and 4.1 cite prior results;
Lemma~4.2 is only sketched; and several probabilistic and concurrent details
are promised for a later full paper.  None of those source boundaries is
upgraded merely because the automata-theoretic architecture is correct.

\subsubsection{B\"uchi's 1962 decision method}

\paragraph{Source identity and boundary.}
The complete article is available in a later Stanford printing of the original
1962 proceedings.  Its printed pp.~1--11 occupy file pp.~12--22 of the
commercial preview; printed p.~11 ends with the references, and file p.~23
begins Julia Robinson's article on printed p.~12.  The original citation is
therefore J.~Richard B\"uchi, ``On a Decision Method in Restricted Second Order
Arithmetic,'' in the proceedings edited by Ernest Nagel, Patrick Suppes, and
Alfred Tarski, Stanford University Press, 1962, pp.~1--11.  No DOI appears in
the original.  The familiar Elsevier and Springer DOIs identify the 1966
series republication and the 1990 collected-works reprint, respectively.  The
public preview supplies a complete research copy of this article but declares
no open-content licence, so the account below is an independent reconstruction
\cite{Buchi1962}.

\paragraph{The interpreted formalism \(SC\).}
B\"uchi's name \(SC\) is retained because replacing it silently by the modern
name S1S would erase the source's own presentation.  It stands for an
interpreted language together with its intended domains, not for an automaton
or a word language.  Individual variables range over \(\mathbb N\), monadic
predicate variables range over all subsets of \(\mathbb N\), and the
nonlogical vocabulary contains \(0\) and unary successor; propositional
connectives and quantifiers of both sorts are available.  If \(t\) denotes a
number, the source writes \(t'\) for its successor
\cite[pp.~1--2]{Buchi1962}.

For a tuple of predicate variables
\(\mathbf i=(i_1,\ldots,i_k)\), put
\[
 \mathbf i(n)=
 \bigl(\mathbf 1_{i_1}(n),\ldots,\mathbf 1_{i_k}(n)\bigr)
 \in\{0,1\}^k.
\]
Here \(\mathbf 1_{i_j}(n)=1\) exactly when \(n\in i_j\), and it is \(0\)
otherwise.
Thus an assignment to \(\mathbf i\) is represented by the infinite word
\[
 \mathbf i(0)\mathbf i(1)\mathbf i(2)\cdots
\]
over the \(2^k\) truth-value states.  This is the source's bridge from monadic
set assignments to infinite words.  The word is the input assignment itself;
no separate run has yet been introduced.  B\"uchi's proposed phrase
``sequential calculus'' describes the use of \(SC\) as a condition language
for sequential circuits and does not name a second formal system.

\paragraph{Formula classes and recurrence.}
The source first organizes formulas into classes \(\Sigma_n\) and \(\Pi_n\).
These symbols name formula classes, not word languages unless a denotation is
explicitly in view.  With propositional matrices \(A,B,C\), its displayed
base classes and recursive clauses are
\[
\begin{aligned}
 \Sigma_1:\quad
 &(\exists\mathbf r)\bigl[
 A[\mathbf r(0)]
 \land(\forall t)B[\mathbf i(t),\mathbf r(t),\mathbf r(t')]
 \land(\exists t)C[\mathbf r(t)]\bigr],\\
 \Pi_1:\quad
 &(\forall\mathbf r)\bigl[
 A[\mathbf r(0)]
 \lor(\exists t)B[\mathbf i(t),\mathbf r(t),\mathbf r(t')]
 \lor(\forall t)C[\mathbf r(t)]\bigr],\\
 \Sigma_{n+1}:\quad
 &(\exists\mathbf r)F(\mathbf i,\mathbf r),
 \qquad F\in\Pi_n,\\
 \Pi_{n+1}:\quad
 &(\forall\mathbf r)F(\mathbf i,\mathbf r),
 \qquad F\in\Sigma_n.
\end{aligned}
\]
The source then defines the infinitely-often quantifier by
\[
 (\exists^\omega t)D(t)
 \quad\Longleftrightarrow\quad
 (\forall x)(\exists t)\,[x<t\land D(t)].
\]
It then singles out \(\Sigma_1^\omega\), whose formulas have the form
\[
 (\exists\mathbf r)\;
 A[\mathbf r(0)]
 \land
 (\forall t)B[\mathbf i(t),\mathbf r(t),\mathbf r(t')]
 \land
 (\exists^\omega t)C[\mathbf r(t)]
\]
where \(A,B,C\) are propositional matrices and \(\mathbf r\) is a finite tuple
of auxiliary monadic predicates.  The notation \(\Sigma_1^\omega\) names a
formula class and, under the paper's declared overload, the relations on
infinite predicate sequences which its formulas define.  It differs from the
ordinary first existential level by requiring the final condition infinitely
often rather than merely once, and the distinction is used in the normal-form
and decision theorems below \cite[pp.~2--3]{Buchi1962}.

There is an exact modern automata comparison, but it is a handbook translation
of the displayed quantifiers rather than a tuple printed in the article.
Suppose \(\mathbf r\) has \(\ell\) coordinates.  For this one formula define
\[
\begin{aligned}
 \Gamma&=\{0,1\}^k,
 &Q&=\{0,1\}^{\ell},\\
 Q_0&=\{q\in Q:A[q]\},
 &\Delta(q,a)&=\{q'\in Q:B[a,q,q']\},\\
 F&=\{q\in Q:C[q]\}.
\end{aligned}
\]
These are eliminable names for the finite data extracted from the displayed
formula: \(\Gamma\) is the input alphabet, \(Q\) the
auxiliary state set, \(Q_0\) the permitted initial states, \(\Delta\) the
successor relation, and \(F\) the recurrent state set.  The existential tuple
\(\mathbf r\) supplies a sequence \(q_0q_1\ldots\), and the formula holds of
the input word exactly when some such sequence begins in \(Q_0\), satisfies
\[
 q_{n+1}\in\Delta(q_n,\mathbf i(n))
 \qquad(n\in\omega),
\]
and visits \(F\) infinitely often.  It is therefore extensionally the
permitted-initial-set variant of the nondeterministic B\"uchi word automaton
defined in Definition~\ref{def:p08-buchi-automaton}; that definition uses one
initial state, whereas \(Q_0\) may contain several.  Nondeterminism is material
because \(B[a,q,q']\) need not determine a unique \(q'\).  The 1962 article
does not package these objects as a five-component tuple, call
\(q_0q_1\ldots\) a run in this notation, or use ``B\"uchi acceptance'' as the
name of the recurrence clause.  Applying the handbook-owned conversion above
with this \(Q_0\) gives the single-start convention used by
Definition~\ref{def:p08-buchi-automaton}; it adds one state and leaves the
recognized language unchanged.

\paragraph{Finite words, congruences, and regular formulas.}
For a finite tuple \(\mathbf r\) of auxiliary monadic predicates,
propositional matrices \(A,B,C\), and endpoints \(x\le y\), the source's class
\(\Sigma^0\) consists of the regular formulas
\[
 R(\mathbf i,x,y)\equiv
 (\exists\mathbf r)\bigl[
 A[\mathbf r(x)]
 \land
 (\forall t)\bigl(x\le t<y\Rightarrow
 B[\mathbf i(t),\mathbf r(t),\mathbf r(t')]\bigr)
 \land
 C[\mathbf r(y)]
 \bigr].
\]
Thus the auxiliary recurrence is initialized at \(x\), constrained exactly on
the finite interval \(x\le t<y\), and tested at \(y\).  The symbol
\(\Sigma^0\) stands first for this formula class.  The truth of
\(R(\mathbf i,x,y)\) depends only on
\[
 \mathbf i(x)\mathbf i(x+1)\cdots\mathbf i(y-1).
\]
Accordingly, \(\Sigma^0\) also denotes the finite-word languages defined by
such formulas.  This is a source-declared formula/denotation overload, not an
identification made without context.

For finite words, a congruence \(\sim\) is an equivalence relation satisfying
\[
 u\sim v
 \quad\Longrightarrow\quad
 uw\sim vw\ \text{ and }\ wu\sim wv
 \qquad\text{for every word }w.
\]
Its rank is its number of equivalence classes.  The source calls a finite-word
language \emph{multiperiodic} exactly when it is a union of classes of one
finite-rank congruence.  This name is needed for Lemmas~6 and~10; it denotes a
finite-word language class, not an ultimately periodic infinite word, and is
not redundant with the later restriction to ultimately periodic predicate
sets \cite[pp.~2--4]{Buchi1962}.

\paragraph{The Sequential Lemma and the finite-word bridge.}
Let \(\mathcal E_0,\ldots,\mathcal E_n\) be a finite partition of the finite
words over the states of one infinite sequence \(\mathbf i\).  Lemma~1, the
Sequential Lemma, gives cut points
\[
 x_1<x_2<x_3<\cdots
\]
and one cell \(\mathcal E_c\) such that
\[
 \mathbf i(x_p)\mathbf i(x_p+1)\cdots\mathbf i(x_q-1)
 \in\mathcal E_c
 \qquad(p<q).
\]
The quantifier \(p<q\) matters: every concatenation of consecutive blocks
between selected cut points has the same color, not only each adjacent block.
The proof colors each two-element set of positions by the cell containing the
intervening word and applies Ramsey's theorem to obtain an infinite
homogeneous set.  This proof is printed and will be used in the complement
argument \cite[p.~3]{Buchi1962}.

Lemma~2 identifies \(\Sigma^0\)-definable finite-word languages with behaviors
of finite-state recursions
\[
 \mathbf r(0)=I,
 \qquad
 \mathbf r(t')=J[\mathbf i(t),\mathbf r(t)],
\]
tested by an endpoint output condition: for a finite-word language
\(\mathcal R\),
\[
 \mathbf i(0)\cdots\mathbf i(x-1)\in\mathcal R
 \quad\Longleftrightarrow\quad
 U[\mathbf r(x)].
\]
The source attributes the difficult
direction to Myhill's subset construction and B\"uchi's 1960 paper; it renders
the converse as a regular formula.  Lemmas~3--5 supply Boolean closure and the
needed bounded and endpoint quantifier reductions.  Lemma~6 states that a
finite-word language is such an automaton behavior exactly when it is
multiperiodic, but refers its proof to Rabin and Scott.  Hence the exact bridge
used here is
\[
 \Sigma^0\text{-definable}
 \quad\Longleftrightarrow\quad
 \text{finite-automaton behavior}
 \quad\Longleftrightarrow\quad
 \text{multiperiodic},
\]
with the cited proof dependencies left visible \cite[pp.~3--4]{Buchi1962}.

\paragraph{Boolean closure and the infinite-word normal form.}
Lemmas~7 and~8 prove closure of \(\Sigma_1^\omega\) under union and
intersection.  The conjunction construction adds finite control which
alternates between waiting for the two recurrent conditions.  Lemma~9 proves
closure under complement.  For a formula with \(m\) auxiliary states, two
finite input words are declared equivalent when, for every ordered pair of
auxiliary states \((V,W)\), they agree both on the existence of a permitted
path from \(V\) to \(W\) and on the existence of one passing through a
designated recurrent state.  This is the intersection of \(2m^2\) binary
distinctions, so its rank is at most
\[
 2^{2m^2}.
\]
The Sequential Lemma supplies a homogeneous infinite block decomposition.
Finitely many initial-block and repeated-block classes can then be separated
according to whether the original formula holds, and the opposite classes
define its complement.  This is an effective finite-congruence complement
argument; it is not a determinization construction silently imported from a
later automata theory \cite[pp.~5--7]{Buchi1962}.

Lemma~10 extracts the structural form
\[
 \mathcal C\,\mathcal D\,\mathcal D\,\mathcal D\cdots,
\]
with \(\mathcal C\) and \(\mathcal D\) multiperiodic finite-word languages.
The display means that an infinite input has some factorization into one
finite stem block from \(\mathcal C\) followed by endlessly many nonempty
finite blocks from \(\mathcal D\).  Every \(\Sigma_1^\omega\)-definable infinite-word
language is a finite union of such sets.

\paragraph{Theorem 1: collapse to one recurrent existential form.}
Lemma~11 states that every \(SC\)-formula can first be put into some level
\(\Sigma_n\); its proof is referred to the method of the 1960 paper rather
than printed here.  A first-level existential formula can be changed into
\(\Sigma_1^\omega\), and repeated use of Lemma~9 handles each alternation.
Theorem~1 therefore states that, for every \(SC\)-formula
\(A(\mathbf i)\), there is a \(B(\mathbf i)\in\Sigma_1^\omega\) equivalent to
it for every assignment to \(\mathbf i\).
The source also says that \(\Sigma_1^\omega\) and \(\Sigma_2\) define the same
relations, so the displayed hierarchy collapses at level two.  This is a
normal-form and definability result; its effective content depends on the
later decision and finite-enumeration argument, not on an unprinted modern
translation bound.  The paper extends the normal form to formulas with free
individual variables and records consequences for the variant in which monadic
predicate assignments are restricted to finite subsets of \(\mathbb N\)
\cite[p.~7]{Buchi1962}.

The finite-data comparison above gives a handbook corollary whose ownership is
important.  For every \(SC\)-formula \(A(\mathbf i)\), the infinite truth-value
words satisfying \(A\) are accepted by a nondeterministic B\"uchi automaton
over \(\{0,1\}^k\).  Conversely, encode the finitely many states of any such
automaton by a finite tuple of predicate bits; propositional matrices can
enforce its initial codes, transitions, and recurrent codes while excluding
unused bit patterns.  The resulting \(\Sigma_1^\omega\)-formula defines exactly
its accepted words.  Theorem~1 supplies the first existence direction, and the
paper's later effectivity argument makes it constructive.  The article itself
does not state this equivalence using the modern automaton tuple or theorem
name.

\paragraph{Lemma 12 and Theorem 2: the decision method.}
After existential predicates are exposed, a sentence in the recurrent normal
form has matrix
\[
 K[\mathbf r(0)]
 \land
 (\forall t)H[\mathbf r(t),\mathbf r(t')]
 \land
 (\exists^\omega t)L[\mathbf r(t)].
\]
If \(\mathbf r\) has \(k\) coordinates, there are \(2^k\) truth-value states.
Any satisfying infinite sequence has an \(L\)-state \(U\) which occurs
infinitely often.  It therefore contains a finite \(H\)-path from a
\(K\)-state to \(U\) and a nonempty finite \(H\)-cycle from \(U\) back to
itself.  Conversely, repeating such a cycle gives a satisfying infinite
sequence.  A path between fixed states need not exceed \(2^k\) edges: a longer
path repeats a state, and deleting the intervening loop preserves the endpoints
and the local \(H\)-constraints.  Only finitely many candidate stems, cycles,
and recurrent states must therefore be considered.

The paper then checks that the earlier normal-form transformations are
effective.  In particular, the congruence used in complementation has an
explicit finite bound, and membership of a candidate pair of homogeneous
classes is reduced to another \(\Sigma_1^\omega\)-sentence decided by
Lemma~12.  Theorem~2 concludes that there is an effective method for deciding
the truth of every \(SC\)-sentence.
The result is decidability, not a modern complexity upper bound.  The finite
enumeration must not be promoted to an elementary, PSPACE, or optimal
algorithm.  The article motivates the result by finite-state circuit
verification but supplies no worked circuit instance \cite[pp.~7--9]{Buchi1962}.

\paragraph{Consequences recorded after Theorem 2.}
The paper then lists five comparisons and consequences, none with a new proof.
First, it says that decidability for what it prints as
\(\Sigma_2^0\)-sentences contains a result of Joyce Friedman and entails
finite-automata algorithms associated with Church and Wang.  This
superscripted notation is reported as the source's formula-class label; the
paper does not define it in that paragraph, and it is not identified here with
the earlier regular-formula class \(\Sigma^0\).  Second, \(SC\) defines
equality, strict order, and congruence modulo each positive integer, which the
source says improves a result of Putnam.

Third, finiteness of a predicate assignment is \(SC\)-definable.  The source
therefore obtains decidability of \(SC_{\mathrm{fin}}\), meaning the same
interpreted formalism with monadic variables restricted to finite subsets of
\(\mathbb N\); it cites B\"uchi's 1960 paper and reports Robinson's attribution
of the result to Ehrenfeucht.  Fourth, taking \(\mathrm{Nn}\) as the natural
numbers and \(\mathrm{Pw}\) as the integral powers of two, it obtains
decidability of the first-order theory
\([\mathrm{Nn},+,\mathrm{Pw}]\), improving Presburger's result.  Fifth, it
relates Theorem~2 to the Skolem--Behmann decision theorem for the monadic
predicate calculus of second order and says a modified Lemma~11 yields a simple
solution.  These are the article's historical comparisons; the eleven pages do
not reprove the cited results \cite[p.~9]{Buchi1962}.

\paragraph{Ultimately periodic assignments and arithmetic consequences.}
A preliminary concluding remark says that a careful analysis of the decision
method would yield a complete axiom system for \(SC\).  It proposes as a
candidate schema the homogeneous-sequence fragment of the Sequential Lemma
used in the complement proof, instantiated for each regular formula
\(E\in\Sigma^0\).  The article neither states nor proves a soundness-and-
completeness theorem for that proposed system.

A set \(X\subseteq\mathbb N\) is ultimately periodic here exactly when
\[
 \exists N\in\mathbb N\;\exists p\in\mathbb N\;
 \bigl[p\ge 1\land
 (\forall n\ge N)\,(n\in X\longleftrightarrow n+p\in X)\bigr].
\]
Let \(SC_{\mathrm{per}}\) be the same interpreted formalism with monadic
variables restricted to those subsets.  This name is necessary to distinguish
a restricted assignment domain from the full monadic domain of \(SC\); it is
defined here, stands for an interpreted formalism, and is used immediately.
Theorem~3 says
\[
 \varphi\text{ is true in }SC_{\mathrm{per}}
 \quad\Longleftrightarrow\quad
 \varphi\text{ is true in }SC
\]
for every \(SC\)-sentence \(\varphi\).  The source outlines the replacement of
Ramsey's theorem by idempotent powers in a finite semigroup rather than giving
the earlier level of detail \cite[pp.~9--10]{Buchi1962}.

On printed p.~10, \(\mathrm{Re}\), \(\mathrm{Ra}\), and \(\mathrm{Pw}\)
name, respectively, the nonnegative reals, the rationals, and the integral
powers of two.  The source's meaning of \(\mathrm{Nn}\) is discrepant: p.~1
calls it the natural numbers, whereas p.~10 calls it the integers.  With those
source abbreviations stated first, Theorem~4 says that the first-order theories
of
\[
 [\mathrm{Re},+,\mathrm{Pw},\mathrm{Nn}]
 \quad\text{and}\quad
 [\mathrm{Ra},+,\mathrm{Pw},\mathrm{Nn}]
\]
are, in its phrase, ``arithmetically equivalent,'' and are decidable.  The
article does not define that phrase or display reductions between the two
theories.  This presentation therefore records it only as an unresolved source
phrase, not as a technical relation used later and not as a choice among
interpretability, reducibility, or definitional equivalence.  The exact
independently usable conclusion is that both displayed first-order theories
are decidable.  The reader likewise does not decide whether the last source
occurrence of ``integers'' was intended to mean nonnegative integers.

The paper attributes undecidability after adjoining \(2x\) to Robinson.  It
also reports, with attribution to G\"odel and without proof, that replacing
monadic predicate quantification by quantification over monadic functions makes
all recursive relations definable.

Three open problems close the article.  Problem~1 calls the formalism with
\(2x+1\) and \(2x+2\) as primitives in place of successor \(SC^2\) and asks
whether it is decidable.  Problem~2 lets \(SC(\alpha)\) mean the variant whose
individual domain is the ordinal \(\alpha\) and whose well-order is primitive,
then asks whether \(SC(\omega^2)\) is decidable.  Problem~3 asks for an
algorithm deciding whether an arbitrary condition
\(C(\mathbf i,\mathbf u)\) is realizable by some finite-state recursion.  The
last is a synthesis question.  Theorem~2 decides whether a supplied
finite-state behavior meets a condition; it does not decide whether some
finite-state implementation exists \cite[p.~10]{Buchi1962}.

\paragraph{Proof and source-strength boundary.}
The article prints the Ramsey proof of Lemma~1, the Boolean constructions and
substantial complement argument of Lemmas~7--9, the finite stem--cycle proof of
Lemma~12, and the remaining effectivity argument for Theorem~2.  Lemma~2's
difficult direction is credited to Myhill and the 1960 paper, Lemma~6 is
referred to Rabin--Scott, Lemma~11 is stated with its method referred
elsewhere, and Theorem~3 receives a finite-semigroup outline.  The modern
finite-data comparison above is handbook-owned but extensionally
exact; the source itself gives no determinization, modern automaton
complementation theorem, translation-state complexity, or circuit case study.

\subsubsection{McNaughton's 1966 presentation}
The publisher identifies the ten-page article and marks its version of record
as open archive, but the article body was not available through the lawful
access routes used here: the publisher PDF route required human verification,
its full-content API required authentication, and the advertised OpenAlex content
copy required a personal API key.  This is an access boundary, not evidence
that the publisher PDF does not exist \cite{McNaughton1966}.

McNaughton's own contemporary Project MAC report supplies one narrower
first-person account.  It explicitly glosses an \(\omega\)-event as a set of
infinite sequences.  It also uses the unanalyzed phrases ``regular
\(\omega\)-event'' and ``finite-state'' and reports, as a Kleene--Myhill
analogue, an equivalence between them \cite[p.~58]{MITProjectMAC1966}.  The
last two phrases are quotations from the report, not predicates defined or
used by the handbook: their exact definitions, their distinction, and the
article's theorem wording remain owed until the article body is read.  The
report likewise does not give the article's alphabet, automaton tuple,
acceptance clause, two constructions, examples, or proof.

The same report places finite-state infinite games in a separately lettered
entry about a 1965 memo and records that the memo's claimed proof had a deep,
then-unpatched fallacy \cite[p.~58]{MITProjectMAC1966}.  No determinacy or
finite-state winning-strategy theorem is therefore imported from that memo
into the 1966 article.  Conversely, no modern deterministic Muller/Rabin tuple,
acceptance family, construction, or size bound is attributed to the article
until its own body is read.

\subsubsection{Safra's 1988 proceedings paper and expanded manuscript}
The registered nine-page proceedings paper remains available here only through
its authoritative abstract and bibliographic record.  The abstract announces a
single-exponential conversion from nondeterministic B\"uchi automata to
deterministic Rabin automata, describes the construction as essentially
optimal, and reports applications to complementation and decision procedures
\cite{Safra1988}.  Safra's university directory, however, also supplies a
different source object: a twenty-page author-hosted PostScript manuscript with
the same title.  Its title-page note calls the 1988 paper a preliminary version,
its own bibliography says that an extended version was to appear in the
\emph{Journal of Computer and System Science}, using that singular printed
wording, and its PostScript header records
generation on 21 March 1995.  No journal publication record was established.
The detailed construction below is therefore attributed only to that expanded
manuscript, not silently to the unread nine proceedings pages
\cite{Safra1995Extended}.

\paragraph{Source automata and pair convention.}
In the expanded manuscript an \(\omega\)-automaton over a finite alphabet
\(\Sigma\) is
\[
  \mathcal A=(\Sigma,Q,q_0,\delta,\mathcal C),
  \qquad
  \delta:Q\times\Sigma\longrightarrow\mathcal P(Q),
\]
where \(Q\) is finite, \(q_0\in Q\), and \(\mathcal C\) is an acceptance
condition.  The input automaton may have \(\delta(q,a)=\varnothing\).  Safra
calls an automaton deterministic exactly when
\(|\delta(q,a)|=1\) for every \(q\in Q\) and \(a\in\Sigma\); this is a total,
single-valued target convention, unlike the partial convention used earlier in
this part.  A run on \(\alpha\in\Sigma^\omega\) is
\(\rho\in Q^\omega\) with \(\rho(0)=q_0\) and
\(\rho(i+1)\in\delta(\rho(i),\alpha(i))\) for every \(i\).  Put
\[
  \operatorname{Inf}(\rho)
  =\{q\in Q:\forall m\;\exists i\ge m\;(\rho(i)=q)\}.
\]
For B\"uchi acceptance, \(\mathcal C=F\subseteq Q\) and a run is accepting
when \(\operatorname{Inf}(\rho)\cap F\ne\varnothing\).  A Muller condition is
a family \(\mathcal M\subseteq\mathcal P(Q)\), and it accepts \(\rho\) exactly
when \(\operatorname{Inf}(\rho)\in\mathcal M\).  In the source's Rabin
pair order, a family \(\{(L_i,U_i):i\in I\}\) accepts a run when, for some
\(i\), a state of \(L_i\) recurs infinitely often and every state of \(U_i\)
occurs only finitely often.  Thus \(L_i\) is the recurring or ``good''
component and \(U_i\) the eventually absent or ``bad'' component.  This order
must not be exchanged with Thomas's later notation merely because both are
called Rabin pairs
\cite[pp.~2--4]{Safra1995Extended}.

\paragraph{The expanded manuscript's determinization state.}
\label{sec:p09-safra-determinization}
Let \(n=|Q|\), fix an \(n\)-element vocabulary \(V\), and fix a linear order on
\(V\) solely to make name allocation deterministic.  A finite ordered
\(V\)-tree is \(T=(N,r,p,S)\), where \(N\subseteq V\), \(r\in N\) is the
root, \(p:N\setminus\{r\}\to N\) is the parent map, and \(S\) linearly orders
the children of each node.  The source defines ancestor strictly; with that
convention, its printed left-to-right clause does not compare siblings because
it asks for strictly ancestral nodes which are siblings.  The handbook owns the
following correction.  For two non-ancestral nodes, follow their root-to-node
paths through their last common node and order them by the first distinct
children, using those children's common-parent sibling order.  For the
handbook-owned deterministic allocation below, extend this corrected relation
to a total traversal order: an ancestor precedes every proper descendant, while
two non-ancestral nodes occur in the corrected left-to-right order.  Unlike the
left-to-right relation, this order also compares ancestors with descendants and
is used only to enumerate surviving temporary nodes.  Each determinization
state printed by the source is exactly
\[
  (T,\ell,c),
  \qquad
  \ell:N\longrightarrow\mathcal P(Q),
  \qquad
  c:V\longrightarrow\{\mathsf{white},\mathsf{green},\mathsf{red}\},
\]
subject to two invariants: the union of the children labels is a proper subset
of the parent label, and labels on non-ancestral nodes are disjoint.  Under the
same strict-ancestor convention, the printed specificity clause, read
literally, would require a state both to belong and not belong to the node
itself.  Its following parenthetical gives the intended inclusive-path
reading.  The handbook therefore makes that correction explicit: the inclusive
root-to-\(d\) path is \(d\) together with all its strict ancestors, and an
original state \(q\) is \emph{specific} to \(d\) exactly when
\(q\in\ell(d)\) and \(q\notin\ell(e)\) for every node \(e\) off that inclusive
path.  Label containment then puts \(q\) in every ancestor label, and branch
disjointness makes \(d\) unique.  ``Specific'' is a metalinguistic
abbreviation for this label relation, not an additional automaton component.
It is needed below to prove that no valid tree has more than \(n\) nodes.  Tree-node
identity is supplied by membership in \(V\); colors do not name nodes, and
state-subset labels do not persist as names
\cite[pp.~9--10]{Safra1995Extended}.

Choose the least \(v_0\in V\) as a handbook-owned representative for the
source's unspecified one-node initial tree.  Initially \(N=\{v_0\}\), the
root label is \(\{q_0\}\), and every vocabulary element is white.  Fixing
\(v_0\), rather than choosing an unnamed root afresh, adds no mathematical
structure: it only makes the displayed deterministic transition a function.

For \(R\subseteq Q\) and \(a\in\Sigma\), define
\[
  \widehat\delta(R,a)=\bigcup_{q\in R}\delta(q,a).
\]
To read \(a\), the expanded manuscript performs the following ordered update.
Here \(V'\) is a temporary \(n\)-element set disjoint from \(V\).

\begin{enumerate}
\item Replace every current label \(R\) by
  \(\widehat\delta(R,a)\).
\item Simultaneously for every node present after the first step whose new
  label \(R\) meets \(F\), attach a youngest child with a fresh temporary name
  from \(V'\) and label \(R\cap F\).
\item Whenever an original state occurs on two non-ancestral paths, retain it
  on the leftmost path and delete it from every path to the right.
\item Delete every empty-labelled node; label containment deletes its
  descendants as well.
\item If a surviving old node's label equals the union of its children labels,
  delete all of that node's descendants and color the node green.  Color every
  other surviving old node white, and color every vocabulary name not retained
  by an old survivor red.
\item Enumerate the surviving temporary nodes in the total traversal order
  defined above and rename them, in that order, by the least unused elements
  of \(V\).
\end{enumerate}

The fixed order in the last step is a handbook-owned choice of representative
for the source's instruction to use ``some'' unused names.  It is necessary
only because the output automaton is asserted to be deterministic.  It does
not identify a new kind of node or change the language.  A newly allocated
name is red in its creation state; if its node survives later updates, the same
name persists and is thereafter white except on a collapse update, when it is
green.  After duplicate removal and collapse, every surviving node has at
least one original state specific to it.  Hence at most \(n\) nodes survive,
so the \(n\) names in \(V\) suffice
\cite[pp.~10--11]{Safra1995Extended}.

There is one definedness case which the printed six steps do not state.  If
advancing the root produces the empty label, step~4 would delete the root even
though every printed state is a rooted tree.  Because the source permits
\(\delta(q,a)=\varnothing\), this case is possible.  The handbook therefore
adds one state \(\bot_{\mathrm{dead}}\), exactly the deterministic automaton
state reached when no original run extends the input prefix.  It is not a tree
node or an additional mathematical structure: it is an eliminable rejecting
sink, loops on every letter, and rejects every continuation.  The next display
states exactly how it enters the acceptance condition.  This addition makes
the transition total and leaves the language and asymptotic bound unchanged.

For each \(v\in V\), put
\[
  G_v=\{(T,\ell,c):c(v)=\mathsf{green}\},
  \qquad
  B_v=\{(T,\ell,c):c(v)=\mathsf{red}\}\cup
      \{\bot_{\mathrm{dead}}\}.
\]
The letters \(G\) and \(B\) are handbook notation for the source's green
\(L_v\) and red \(U_v\) components.  They are used only to prevent a collision
with language notation.  The added state \(\bot_{\mathrm{dead}}\) lies in every
\(B_v\) and no \(G_v\).  The deterministic run is accepting exactly when
some \(v\) is red only finitely often and green infinitely often, equivalently
when its infinity set meets \(G_v\) and misses \(B_v\) for some \(v\).  The
construction therefore has exactly \(n\) Rabin pairs in this expanded
presentation
\cite[p.~11]{Safra1995Extended}.

\paragraph{Why the tree and the colors suffice.}
The source first shows why a plain subset state loses lineage information and
why a single root-level check that all currently represented states have
visited \(F\) since the preceding successful check can still be incomplete.
The tree applies its already defined child-creation and collapse operations
recursively below each node: a child tracks the accepting-state mass seen
within its parent's currently represented runs, horizontal deletion gives
responsibility for a merged original state to the older branch, green records
equality of a node label with the union of its child labels at a collapse, and
red records loss of the named lineage.  In the manuscript's worked overview,
one accepted word is witnessed
by infinitely many root collapses, while another is witnessed by a stable leaf
that collapses infinitely often.  The examples demonstrate the two ways the
Rabin condition can select a persistent node; they are not a definition of the
construction
\cite[pp.~5--9]{Safra1995Extended}.

Fix an input word \(\alpha\in\Sigma^\omega\) and the unique run of the
determinized automaton on \(\alpha\).  For \(i\le j\), write
\(\alpha[i,j)\) for the finite factor
\(\alpha(i)\alpha(i+1)\cdots\alpha(j-1)\).  For soundness, suppose one name
\(v\) in that determinized run is eventually never red and is green at
positions
\[
  i_1<i_2<\cdots.
\]
Let \(S_j\) be the label of \(v\) at \(i_j\), and take
\(S_0=\{q_0\}\), \(i_0=0\).  The child-creation and collapse clauses imply
that every \(q\in S_{j+1}\) has some predecessor \(q'\in S_j\) joined to it
by an \(\mathcal A\)-run over
\(\alpha[i_j,i_{j+1})\) which visits \(F\).  Form the finitely branching tree
whose level \(j\) consists of \((q,j)\) for \(q\in S_j\), choosing such a
predecessor for every next-level node.  K\"onig's lemma supplies an infinite
branch.  Concatenating its finite run segments gives one original run which
visits \(F\) between every two successive green positions, hence is B\"uchi
accepting
\cite[pp.~11--12]{Safra1995Extended}.

For completeness, let \(\rho\) be an accepting run of the original automaton
\(\mathcal A\) on the same word \(\alpha\).  Its current state is always in
the root label.  If the root turns green infinitely often,
its name is the required Rabin witness.  Otherwise, after the root's last green
event, the next visit of \(\rho\) to \(F\) places its state in a child.  It may
move left to an older sibling only finitely many times, so it eventually stays
under one surviving child.  Either that child turns green infinitely often, or
the same argument descends another level.  Proper child containment bounds the
depth by \(n\); the descent must therefore stop at a surviving node whose name
is eventually never red and is green infinitely often
\cite[p.~12]{Safra1995Extended}.

The number of ordered trees on subsets of \(V\) is at most
\(8^n n!=2^{n\log n+O(n)}\).  Indeed, there are \(2^n\) choices of the
node set.  For a set of \(k\le n\) nodes, an ordered rooted tree shape is
encoded by the balanced-parenthesis word of its depth-first traversal, so
there are at most \(2^{2k}\le4^n\) shapes; assigning the chosen names to its
nodes gives at most \(k!\le n!\) further choices.

For a fixed tree, the nodes whose labels contain a given original state
\(q\) form either the empty set or one inclusive root-to-node path:
label containment makes the set upward closed toward the root, and
disjointness of non-ancestral labels prevents branching.  Thus a present
state has a unique deepest, specific node.  Fix \(\star\notin V\) and encode
the labeling by
\[
 f:Q\longrightarrow N\cup\{\star\},
\]
sending a present state to its specific node and an absent state to
\(\star\).  This encoding determines every label: \(q\in\ell(d)\) exactly
when \(f(q)\in N\) and \(d\) lies on the inclusive root-to-\(f(q)\) path.
It is therefore injective on valid labelings.  There are at most
\((|N|+1)^n\le(n+1)^n\) encodings, whether or not every such function
gives a valid labeling.  The symbol \(\star\) is only an encoding value.
The color map has at most \(3^n\) possibilities, and the rejecting sink
contributes one further automaton state.  Hence the total is at most
\[
 1+8^n n!\,(n+1)^n3^n=2^{O(n\log n)}.
\]
Together with the correctness proof, this establishes Theorem~1: every
\(n\)-state nondeterministic B\"uchi automaton has an equivalent deterministic
Rabin automaton with \(2^{O(n\log n)}\) states and exactly \(n\) Rabin pairs
\cite[p.~12]{Safra1995Extended}.
The adjoining discussion cites a modified Michel argument for an \(n!\)
lower bound and therefore calls the construction optimal only up to a constant
in the exponent.  Lemma~2 separately says that conversion to deterministic
Muller acceptance can require a family \(\mathcal M\) with doubly exponentially
many members; its proof is only sketched in Appendix~A
\cite[pp.~12--13, 17--18]{Safra1995Extended}.

The applications preserve their own target classes.  Corollary~4 complements
an \(n\)-state nondeterministic B\"uchi automaton by a deterministic Rabin
automaton with \(2^{O(n\log n)}\) states and \(n+1\) pairs.  Corollary~6 obtains
a complementary nondeterministic B\"uchi automaton with
\(2^{O(n\log n)}\) states.  Safra calls a B\"uchi automaton
\emph{deterministic in the limit} exactly when every accepting run has only
finitely many positions at which its current state and input letter admit more
than one successor.  This is a property of the automaton's accepting runs and
is weaker than full determinism: rejecting runs need not eventually cease
making nondeterministic choices.  Corollary~7 gives this different, partial
determinization, producing an equivalent B\"uchi automaton deterministic in the
limit with \(O(3^n)\) states.  The source then records the resulting
exponential and doubly exponential upper bounds for the modal-logical decision
procedures it discusses.  These are consequences in the expanded manuscript, not extra
claims recovered from the unread proceedings body
\cite[pp.~13--16]{Safra1995Extended}.

This source distinction is mathematically material.  Safra's expanded
manuscript uses an \(n\)-element persistent vocabulary and obtains \(n\) pairs;
Thomas's 1996 report presentation below uses names from
\(\{1,\ldots,2n\}\) and prints \(2n\) pairs.  The handbook keeps both
constructions under their actual sources.  It does not combine the smaller
pair count from one with the update conventions of the other, and it does not
claim that the author-hosted expansion is a facsimile of the 1988 paper or a
published journal article.

\subsubsection{Thomas's 1996 report precursor}
Thomas's complete May 1996 Kiel report precursor supplies the later survey
presentation at Sections~5.1--5.3.  Renaming only the source alphabet and
transition relation to avoid collisions, Definition~5.1 has the form
\[
 \mathcal A=(Q,\Sigma,q_0,\Delta,\operatorname{Acc}),
 \qquad
 \Delta\subseteq Q\times\Sigma\times Q,
\]
where \(Q\) and \(\Sigma\) are finite and \(q_0\in Q\).  A run of \(\mathcal A\)
on \(\alpha\in\Sigma^\omega\) is a sequence \(r\in Q^\omega\) satisfying
\[
 r(0)=q_0,
 \qquad
 (r(i),\alpha(i),r(i+1))\in\Delta
 \quad(i\in\mathbb N),
\]
and the source writes
\[
 \operatorname{In}(r)
 =\{q\in Q:(\forall m\in\mathbb N)(\exists i\ge m)\ r(i)=q\}.
\]
This is the exact eliminable meaning of the source's shorthand
\(\exists^\infty i\;r(i)=q\): the state occurs at arbitrarily late run
positions.  The set stands for the recurrent states of this run and is used in
each acceptance clause below.
When the automaton is deterministic, Thomas replaces \(\Delta\) by a total
transition function
\[
 \delta:Q\times\Sigma\longrightarrow Q,
 \qquad
 r(i+1)=\delta(r(i),\alpha(i))
 \quad(i\in\mathbb N).
\]
Thus every input word has one run.  Totality, rather than the weaker condition
that there be at most one successor, is the premise used below when the Muller
family is complemented.
The acceptance component has one of four explicitly different standings.  For
\(F\subseteq Q\), B\"uchi acceptance is
\(\operatorname{In}(r)\cap F\neq\varnothing\).  For
\(\mathcal F\subseteq\mathcal P(Q)\), Muller acceptance is
\(\operatorname{In}(r)\in\mathcal F\).  For a finite sequence
\(((E_i,F_i))_{1\leq i\leq m}\), Rabin acceptance is
\[
 \exists i\leq m\;
 \bigl(\operatorname{In}(r)\cap E_i=\varnothing
       \land
       \operatorname{In}(r)\cap F_i\neq\varnothing\bigr),
\]
whereas Thomas's Streett convention is its dual
\[
 \forall i\leq m\;
 \bigl(\operatorname{In}(r)\cap E_i\neq\varnothing
       \lor
       \operatorname{In}(r)\cap F_i=\varnothing\bigr).
\]
Thus, with the source's ordering of a Streett pair, recurrence of some
\(F_i\)-state requires recurrence of some \(E_i\)-state.  Example~5.2 presents
the language ``if \(a\) occurs infinitely often, then so does \(b\)'' in all
four forms, and Proposition~5.3 proves that nondeterministic automata with any
of the four modes recognize the same \(\omega\)-languages.  The comparison is
therefore source-carried rather than supplied from modern terminology
\cite[\S5.1]{Thomas1996Report}.

Thomas also records the exact finite-union characterization of this class.  An
\(\omega\)-language \(L\subseteq\Sigma^\omega\) is B\"uchi-recognizable if and
only if, for some \(k<\omega\), there are regular finite-word languages
\(U_j,V_j\subseteq\Sigma^*\) such that
\[
 L=\bigcup_{j<k}U_j\cdot V_j^\omega.
\]
Here \(V^\omega\) is the set of infinite concatenations
\(v_0v_1\cdots\in\Sigma^\omega\), with every \(v_i\in V\), for which
\(\lim_{m\to\infty}|v_0\cdots v_m|=\infty\).  Thus empty factors cannot turn an
infinite concatenation into a finite word and may equivalently be deleted.
This is Thomas's report-level survey statement, not original-priority evidence
for the underlying theorem \cite[p.~30]{Thomas1996Report}.

Determinism changes the boundary.  A deterministic B\"uchi automaton recognizes
the words with infinitely many \(a\)'s, but no deterministic B\"uchi automaton
recognizes their complement; deterministic Muller acceptance is closed under
complement by replacing \(\mathcal F\) with
\(\mathcal P(Q)\setminus\mathcal F\).  Theorem~5.4 states McNaughton's
effective conversion from a nondeterministic B\"uchi automaton to an equivalent
deterministic Muller automaton and proves it by Safra's construction
\cite[\S5.2]{Thomas1996Report}.

The construction needs more information than a subset of reachable states,
because recurring membership in successive subsets need not belong to one
infinite accepting run.  Put \(n=|Q|\).  A \emph{Safra tree} in this
presentation is a tuple \((T,\ell,\nu,M)\), where \(T\) is a finite rooted
ordered tree with live-node set \(N(T)\),
\(\ell:N(T)\to\mathcal P(Q)\) labels its nodes by original-state subsets,
\(\nu:N(T)\hookrightarrow\{1,\ldots,2n\}\) assigns distinct names to distinct
live nodes, and \(M\subseteq N(T)\) is the set of nodes carrying the source's
breakpoint mark.  Sibling labels are disjoint, and the union of the child
labels is a proper subset of their parent's label.  This tuple is a state of
the determinized word automaton, not an input tree and not a run of Rabin's
tree automaton.  Its name map is distinct from both the changing state-subset
labels and the current-update marks: it alone identifies a surviving lineage,
which is why persistence and the Rabin pairs below require it.  The initial
tree has one root, labelled \(\{q_0\}\), named \(1\), and unmarked.  Fixing
the initial name as \(1\) is a handbook-owned normalization; any other single
initial name differs only by a renaming.  A mark records a collapse in the
immediately preceding letter update; it is not a persistent property of a
node.  To make that event meaning literal, the
handbook clears every old mark before the next update while retaining the node
names.  Only a collapse during the current update marks the successor state.
This mark-lifetime clause is a handbook-owned clarification of the printed
transition.  The source then performs the following operations in order:
\begin{enumerate}
\item for every label \(R\) meeting the B\"uchi set \(F\), add \(R\cap F\) as
  its youngest child;
\item advance every label, including those just created, by the transition
  relation;
\item whenever a state in a younger sibling's label also occurs in an older
  sibling's label, delete it from the younger sibling and from every label in
  that younger sibling's subtree; remove every empty-labelled node together
  with its descendants; then collapse all children of a parent whose
  child-label union equals its label, marking that parent.
\end{enumerate}
Deleting through the younger subtree is required to preserve containment of
each descendant label in its parent label.  This is a handbook-owned
definedness clarification of the source's compressed horizontal-deletion
instruction, not an additional acceptance condition.
The removal clause applies to a nonempty advanced root.  If the root image in
the second stage is empty, the successor is instead the empty ordered tree:
the determinized word-automaton state having no node, name, label, or mark.
Every later letter leaves this state unchanged.  This additional state is
necessary because deleting an empty root would not leave a rooted tree, and it
differs from every ordinary construction state precisely by having no root.  If
the advanced root is empty, no original run extends the prefix just read, and
no continuation can restore such a run; the absorbing empty state therefore
preserves rejection.  It has zero live nodes, uses no node name, and adds no
accepting pair, so it leaves the node, name, state-asymptotic, and pair bounds
below unchanged.  This is a handbook-owned completion of the printed update.
Names of surviving old nodes persist.  Fix the natural order
\(1<\cdots<2n\).  After cleanup, list the surviving newly created nodes by the
depth-first order which lists a node before its descendants and an older
sibling's subtree before a younger sibling's subtree.  In that order assign
each new node the least name which occurs neither in the entire predecessor
name set nor on an earlier new node in the same update.  This handbook-owned
choice makes the successor name map unique, prevents a deleted lineage from
being confused with a new one, and preserves injectivity.  At most \(n\) nodes
are live, and the names \(1,\ldots,2n\) suffice: at most \(n\) predecessor
names are withheld while at most \(n\) new nodes need names, and cleanup again
leaves at most \(n\) live nodes.  The Muller condition accepts exactly when some name is
eventually permanent and marked infinitely often.  Equivalently, for every
\(k\in\{1,\ldots,2n\}\), let
\[
 E_k=\{(T,\ell,\nu,M):k\notin\operatorname{im}(\nu)\},
 \qquad
 F_k=\{(T,\ell,\nu,M):\exists v\in M\;\nu(v)=k\},
\]
with the empty absorbing state included in every \(E_k\) and no \(F_k\).  The
resulting deterministic Rabin automaton uses the pairs \((E_k,F_k)\).  The
proof derives an original accepting
run from infinitely many current collapse events on a persistent lineage and
proves the converse by descending through the bounded tree height.  The empty
state lies in every \(E_k\) and no \(F_k\), hence is rejecting.  Clearing old
marks is essential: for
\(q_0\mathrel{\xrightarrow{a}}f\mathrel{\xrightarrow{b}}s\) with a
\(b\)-loop at \(s\) and B\"uchi set \(\{f\}\), the rejected word
\(ab^\omega\) produces only one root collapse; retaining that one mark would
falsely place the persistent root in \(F_k\) forever.  Current-update marking
rejects it.  Proposition~5.6 gives
\(2^{O(n\log n)}\) determinization states and exactly \(2n\) pairs in this
presentation.  Theorem~5.7's lower-bound qualification rules out a general
\(2^{O(n)}\)-state, \(O(n)\)-pair conversion, using Theorem~5.8's factorial
B\"uchi-complementation family; it is not an exact-constant claim
\cite[\S5.2]{Thomas1996Report}.

Theorem~5.9 then proves the effective equivalence between B\"uchi-recognizable
\(\omega\)-languages and monadic-second-order definability.  State-set
variables describe an accepting run in one direction; atomic automata,
Boolean closure (using determinization for complement), and track projection
give the other.  Corollary~5.10 derives decidability of \(S1S\) from an
input-free B\"uchi automaton \(\mathcal A\): the sentence is true exactly
when \(\mathcal A\) has an accepting run.  In its finite transition graph,
the exact criterion is that some \(q\in F\) is reachable from the initial
state and has a positive-length path back to itself.  The reaching path
may have length zero; the return path must contain at least one edge.
For necessity, an accepting run visits \(F\) infinitely often.  Since
\(F\) is finite, some \(q\in F\) occurs infinitely often; its first
occurrence gives a reaching path, and two successive occurrences give
a positive-length return.  Conversely, concatenate a reaching path with
infinitely many repetitions of a positive-length return from an accepting
state \(q\) to itself.  This is an infinite run visiting \(q\in F\)
infinitely often, hence accepting.  Both finite graph conditions are
decidable by reachability search.  In particular, an isolated accepting
state without a positive return does not suffice
\cite[Cor.~5.10, p.~38]{Thomas1996Report}.

Proposition~5.11 writes every MSO formula
as a finite disjunction of clauses
\[
 \bigl(\exists x\,\forall y>x\;\neg\varphi_i(\bar X,y)\bigr)
 \land
 \bigl(\forall x\,\exists y>x\;\psi_i(\bar X,y)\bigr),
\]
where every quantifier inside \(\varphi_i\) and \(\psi_i\) is restricted below
\(y\).  Corollary~5.12 concludes that every MSO formula over \(\omega\)-words
is equivalent to a weak-MSO formula, meaning one whose monadic quantifiers
range only over finite sets
\cite[\S5.3]{Thomas1996Report}.

Thomas's report is an authoritative later proof source at explicit precursor
strength.  The registered 1997 chapter is not part of the source basis used
here, so these locators belong only to the report.  They establish nothing
about the literal bodies of other historical works: B\"uchi 1962 is supported
separately by the complete article body treated above, while McNaughton 1966
and Safra 1988 remain unread.  The report also does not convert the
Vardi--Wolper product into an original determinization proof.
Any admission of this text is therefore admission only of the bounded
presentation actually printed here.  It does not close the separate
registered-body obligations for McNaughton 1966, Safra 1988, or Thomas 1997,
and it does not turn the Vardi--Wolper author preprint into a proceedings
facsimile.
